% =====================================================================
%  Partie 4 (~4 min) : mon travail
% =====================================================================
\partie{Mon travail}{Mon travail : des modèles plus humains ?}

% ---------------------------------------------------------------------
\begin{frame}{La question}
  \centering
  \vspace{0.2em}
  {\large Si l'on donne à un modèle des \jaun{limites humaines},\\devient-il plus proche de l'humain ?}

  \vspace{1.0em}
  \begin{columns}[T,onlytextwidth]
    \column{.48\textwidth}
    \visible<2->{\carte[orange]{Contrainte de mémoire}{%
      Le modèle oublie progressivement le contexte, comme notre mémoire de travail.}}
    \column{.48\textwidth}
    \visible<3->{\carte[orange]{Contrainte de données}{%
      10 à 100 millions de mots, ce qu'entend un enfant (défi \textbf{BabyLM}), au lieu de milliers de milliards.}}
  \end{columns}

  \vspace{0.8em}
  \visible<4->{%
  \begin{tikzpicture}
    \foreach \n/\l/\x in {118/modèles/0, 18/jeux de données/3.6, 7/comportement/7.2, 11/cerveau/10.8}{
      \node[font=\bfseries\Large, text=jaune] at (\x,0) {\n};
      \node[font=\scriptsize, text=gris] at (\x,-0.5) {\l};}
  \end{tikzpicture}\\[0.1em]
  {\scriptsize\color{gris}temps de lecture, oculométrie ; IRMf, MEG, EEG, ECoG}}
\end{frame}

% ---------------------------------------------------------------------
\begin{frame}{Résultat : comportement et cerveau ne vont pas ensemble}
  \begin{columns}[c,onlytextwidth]
    \column{.50\textwidth}
    \centering
    \photo{height=4.9cm}{images/resultat_pareto.png}
    \column{.48\textwidth}
    \small
    Chaque point est un modèle. \hum{À droite} : proche du cerveau. \jaun{En haut} : proche de la lecture humaine.
    \begin{itemize}
      \item<2-> Les modèles \textbf{contraints} (\textcolor{violet!60!white}{violet}, \textcolor{rouge}{rouge}) sont plus proches de la \textbf{lecture}\dots
      \item<3-> \dots mais plus \textbf{loin du cerveau}.
      \item<4-> Les grands modèles (\textcolor{bleu}{bleu}) montrent l'inverse.
      \item<5-> \textbf{Aucun} modèle ne gagne sur les deux à la fois.
    \end{itemize}
  \end{columns}
  \vspace{0.3em}
  \visible<6->{\cle{« Ressembler à l'humain » n'est pas une seule propriété : il faut préciser à quel niveau.}}
  \src{Sow et al., travail en cours (LSCP)}
\end{frame}
